\uuid{a6IO}
\titre{Module/argument}
\auteur{bodin ; hanani ; mzari}
\organisation{exo7}
\licence{CC BY-NC-SA 4.0}
\datecreate{2020-01-20}
\niveau{L1}
\module{Algèbre}
\chapitre{Nombres complexes}
\sousChapitre{Forme cartésienne, forme polaire}
\difficulte{3}
\theme{Écritures algébrique et géométrique, module, argument}
\exo7qcmid{lille-1.nombres-complexes.013}

\contenu{
\begin{question}
\qclassification{104.01}
\qtags{motcle=module/argument}
Soit $z=1+ e^{i\theta},\theta \in ]-\pi,\pi[$. Quelles sont les assertions vraies ?
\begin{answers}
    \bad{$|z|=2 $}

    \good{$|z|=2\cos(\frac{\theta}{2}) $}

    \good{$\arg z =  \frac{\theta}{2} \, [2\pi]$}

    \bad{$\arg z = \theta \, [2\pi]$}

\end{answers}
\begin{explanations}
$z=e^{i\frac{\theta}{2}} (e^{i\frac{\theta}{2}} + e^{-i\frac{\theta}{2}}) = 2 \cos (\frac{\theta}{2}) e^{i\frac{\theta}{2}}$. Comme $\theta \in ]-\pi,\pi[$ ,  $\cos (\frac{\theta}{2})>0$. On déduit que : $|z|=2\cos (\frac{\theta}{2})$ et $\arg z =  \frac{\theta}{2} \, [2\pi]$.
\end{explanations}
\end{question}
}
